\Subhead{2.2.3}{Временная дисперсия. Соотношения Крамерса--Кронига и принцип причинности}

В программе фраза «Поле в средах с временной» обрывается. Здесь разобрана естественная для соседних пунктов тема \emph{временной дисперсии}; это пояснение, а не исправленная официальная цитата.

Среда с памятью связывает поляризацию с предысторией поля:
\begin{equation}
 P(t)=\varepsilon_0\int_0^\infty\chi(\tau)E(t-\tau)\dd\tau,\quad
 \chi(\omega)=\int_0^\infty\chi(\tau)\e^{\ii\omega\tau}\dd\tau.
\end{equation}
Нижний предел выражает причинность. При достаточном убывании отклика $\chi(\omega)$ аналитична в верхней полуплоскости комплексной частоты. Формула Коши на контуре, замкнутом там, связывает действительную и мнимую части:
\begin{align}
 \chi'(\omega)&=\frac1\pi\mathcal P\int_{-\infty}^{\infty}
 \frac{\chi''(\omega')}{\omega'-\omega}\dd\omega',\\
 \chi''(\omega)&=-\frac1\pi\mathcal P\int_{-\infty}^{\infty}
 \frac{\chi'(\omega')}{\omega'-\omega}\dd\omega'.
\end{align}
Если отклик не стремится к нулю, вычитают высокочастотную константу. Для вещественной импульсной характеристики $\chi'(-\omega)=\chi'(\omega)$, $\chi''(-\omega)=-\chi''(\omega)$. Поэтому дисперсия и поглощение не являются независимыми функциями. Причинность запрещает передачу новой информации раньше фронта даже там, где формальная групповая скорость больше $c$ или отрицательна.
